% =====================================================================
%  ch3c_relations.tex — §3.3 A∞-关系
% =====================================================================
\subsection{$A_\infty$-关系}\label{ss:relations}

现在陈述并证明：

\begin{theo}\label{thm:ainfty}
上面定义的各操作 $\mu^k$ 满足 $A_\infty$-关系
\begin{equation}\label{eq:ainfty}
\sum_{\ell=1}^{k}\sum_{j=0}^{k-\ell} (-1)^* \mu^{k+1-\ell}\bigl(\rho_k,\dots,
\rho_{j+\ell+1},\mu^\ell(\rho_{j+\ell},\dots,\rho_{j+1}),\rho_j,\dots,\rho_1\bigr)=0,
\end{equation}
其中 $*=j+\deg(p_1)+\dots+\deg(p_j)$。
\end{theo}

证明依据与通常情形相同的几何思想：说明一维模空间给出 \eqref{eq:ainfty} 左端所出现的那些刚性构形之间的配边。但论证需要作实质性的修改，以计入通过 $M$ 的节点的传播。我们从下面的引理开始，它是通常圆盘相应结论的直接推论：

\begin{lemma}\label{l:localslit}
传播全纯圆盘模空间的一维层对应于这样的构形：其中一个圆盘分量带单个边界分支点（除此之外是浸入的），或者是一个带四个标记点（其中一个是输出）的常值映射，而其余所有分量都是刚性的（即带局部凸角的浸入多边形）。
\end{lemma}

\begin{proof}
该结论来自黎曼曲面上全纯圆盘的自动正则性（例如见 \cite[\S13]{SeBook}），以及每个（非常值）圆盘模空间的维数等于边界分支点数加上两倍内部分支点数这一事实，其局部坐标由分支点的像的位置给出（例如见 \cite[Proposition 7.8]{ENS}）。
\end{proof}

带边界分支点的圆盘附近，边界沿一条「狭缝」（slit）折回自身，变形通过沿边界移动分支点来推进，要么把狭缝更深地伸入圆盘，要么把它收缩。这种一维层的端点出现于：
\begin{enumerate}
\item[(i)] 狭缝一直收缩到一个标记点；或
\item[(ii)] 狭缝一直横穿圆盘，把它裂成一对在某个 Floer 生成元或某个节点处相接的圆盘。
\end{enumerate}

我们不再讨论带四个标记点的常值分量，因为唯一需要显式处理它的情形是一对位于节点处的常值圆盘，那并不涉及任何传播。

考察带边界分支点的传播全纯圆盘的整体，折回自身的那部分边界（狭缝）可以延伸到承载分支点的分量之外。也就是说，设传播全纯圆盘 $u:S\to M$ 在映到 $L_i$ 的某条边上有一个边界分支点。那么 $\partial S$ 上从 $z_{i-1}$ 到 $z_i$ 的那一段映到沿 $L_i$ 从 $p_{i-1}$ 走到 $p_i$ 的弧，途中可能经过 $M$ 的节点，并在分支点处恰好折回一次。被走过两次的那部分边界（\emph{狭缝}）可以穿过 $M$ 的节点，此时我们说狭缝\emph{传播}通过这些节点。狭缝的端点要么是 $p_{i-1}$ 或 $p_i$ 之一，此时映射的像有一个\emph{凹角}（若 $p_{i-1}=p_i$ 则称为\emph{重复角}）；要么是 $M$ 的某个节点，从分支点到 $p_{i-1}$ 与到 $p_i$ 的弧在那里分叉进入不同分量，此时我们说映射的像有一个\emph{分叉节点}。这些可能性互斥。

为证明定理~\ref{thm:ainfty}，我们考虑如下一维模空间的并：边界在给定拉格朗日 $L_0,\dots,L_k$ 上、角在固定的输入与输出生成元 $p_1,\dots,p_k,p_0$ 处、且其像 $u(S)$ 以固定重数覆盖 $M$ 中某个固定区域的传播全纯圆盘。如上所述，这个区域必须恰有一个凹角（或重复角）或一个分叉节点。不同的一维层对应于狭缝从该凹角（或重复角）、或从分叉节点出发，沿途经过节点传播，最终抵达某个分量内分支点的不同方式。

我们在端点处把这些模空间互相粘合——只要端点对应于狭缝走过完全相同路径（终止于 $M$ 的某个节点）的传播圆盘——从而得到一维胞腔复形，其端点对应于破裂构形；例如见图~\ref{fig:ainfty}。困难在于：这个胞腔复形一般具有\emph{树}的结构，边朝多个方向分叉，而不是一维流形。树的根对应狭缝收缩成一点的情形；离开根的不同路径对应狭缝可以增长的不同路径。

于是，定理~\ref{thm:ainfty} 证明的关键步骤是我们称为「总传播重数不变性」的性质：狭缝越过节点时，总传播重数保持不变。

\begin{lemma}\label{l:tot_mult_invariance}
一个狭缝恰好停在 $M$ 的某个节点之前的构形的传播重数，等于把狭缝略微延伸越过该节点所能得到的各构形的传播重数之和，再加上狭缝恰终止于该节点处的那些破裂构形的传播重数。
\end{lemma}

在 §\ref{sss:concave} 中我们将说明，对像具有凹角或重复角的构形，该引理给出所需的一维模空间之并的端点之间的抵消；引理本身在 §\ref{sss:proof_lemma_tot} 证明。分叉节点的情形还需要额外的要素，见 §\ref{sss:bifurcated}。

\subsubsection{带一个凹角的传播圆盘}\label{sss:concave}

先考虑一族的传播扰动全纯圆盘，它们有 $k+1$ 个标记点映到 Floer 生成元 $p_0,p_1,\dots,p_k$，且像在其中一个标记点处有凹角，比如 $p_i\in\cX(L_{i-1},L_i)$。这族中最简单的层由狭缝不传播的构形组成；此时 $u(S)$ 的某个分量映到在 $p_i$ 处有凹角的浸入多边形，其余分量皆刚性。这类构形通过沿 $L_{i-1}$ 或沿 $L_i$ 移动边界分支点来变形，从而在多边形中产生一条狭缝，它从凹角出发沿其中一条拉格朗日延伸，如图~\ref{fig:ainfty} 中部所示。

在黎曼曲面上通常的拉格朗日 Floer 理论中，一维模空间的每个连通分量都是区间，其端点出现在狭缝一直横穿、最终碰到凹多边形的边界、把它裂成一对较小的凸多边形之时。这些破裂构形贡献 \eqref{eq:ainfty} 左端 $p_0$ 的系数，而 $A_\infty$-关系表达的正是它们成对抵消。（破裂构形的面积权重与和乐权重与其作为极值变形的未破裂构形相同，因而在区间两端相等。）

\begin{figure}[t]
\centering
\begin{tikzpicture}[scale=0.62]
% (a) 带狭缝的多边形
\begin{scope}[shift={(-4.6,0)}]
\fill[orange!18] (-1.9,0) -- (-0.6,1.2) -- (1.6,1.2) -- (2.4,0) -- (1.6,-1.2) -- (-0.6,-1.2) -- cycle;
\draw[thick] (-1.9,0) -- (-0.6,1.2) -- (1.6,1.2) -- (2.4,0) -- (1.6,-1.2) -- (-0.6,-1.2) -- cycle;
\draw[blue,thick] (-1.9,0) -- (-0.7,0.02);
\draw[blue,thick] (-0.7,0.02) -- (-1.9,0.06);
\node[blue] at (-1.35,-0.35) {\small $L_i$};
\fill (-1.9,0) circle (0.07) node[below left] {\small $p_i$};
\node at (0.2,0.55) {\small $u(S)$};
\node at (0.2,-1.6) {\small (a) 狭缝从凹角出发};
\end{scope}
\draw[->,thick] (-2.3,0)--(-1.5,0);
% (b) 狭缝延伸
\begin{scope}[shift={(0,0)}]
\fill[orange!18] (-1.9,0) -- (-0.6,1.2) -- (1.6,1.2) -- (2.4,0) -- (1.6,-1.2) -- (-0.6,-1.2) -- cycle;
\draw[thick] (-1.9,0) -- (-0.6,1.2) -- (1.6,1.2) -- (2.4,0) -- (1.6,-1.2) -- (-0.6,-1.2) -- cycle;
\draw[blue,thick] (-1.9,0) -- (1.2,0.55);
\draw[blue,thick] (1.2,0.55) -- (-1.9,0.06);
\fill (-1.9,0) circle (0.07);
\node at (0.2,-1.6) {\small (b) 狭缝加长};
\end{scope}
\draw[->,thick] (2.7,0)--(3.5,0);
% (c) 破裂
\begin{scope}[shift={(5.2,0.7)}]
\fill[orange!18] (-1.7,0) -- (-0.3,1.0) -- (1.3,1.0) -- (1.9,0) -- (-1.7,0);
\draw[thick] (-1.7,0) -- (-0.3,1.0) -- (1.3,1.0) -- (1.9,0) -- (-1.7,0);
\fill (1.9,0) circle (0.07) node[right] {\small $p_0$};
\end{scope}
\begin{scope}[shift={(5.2,-0.9)}]
\fill[orange!18] (-1.7,0) -- (1.9,0) -- (1.3,-1.0) -- (-0.3,-1.0) -- (-1.7,0);
\draw[thick] (-1.7,0) -- (1.9,0) -- (1.3,-1.0) -- (-0.3,-1.0) -- (-1.7,0);
\fill (-1.7,0) circle (0.07) node[left] {\small $p_i'$};
\end{scope}
\node at (5.2,-2.3) {\small (c) 狭缝横穿、破裂成两片};
\end{tikzpicture}
\caption{带凹角的传播圆盘的一维族（示意图；原书图 4）。狭缝沿其中一个方向的拉格朗日延伸，最终把多边形裂成一对刚性构形。}\label{fig:ainfty}
\end{figure}

在我们的设定中，当狭缝横穿凹多边形时，可能先碰到 $M$ 的某个节点——圆盘 $u:S\to M$ 经该节点传播——而不是圆盘的边界。此时变形自然继续下去（进入其他模空间），因为可以允许狭缝传播到 $u(S)$ 的下一个分量，如此继续，直到最终碰到传播圆盘的边界。然而，若 $u$ 在局部把狭缝所要伸入的 $M$ 分量覆盖多于一次，则沿相应拉格朗日切开的方式可能不止一种。图~\ref{fig:ainfty} 左侧示意了这一情形：$S$ 最左的分量（以 $L_0$ 与 $L_1$ 为边界的带）以输入次数 $k_{in}=2$ 进入左侧节点，于是这个全纯圆盘沿 $L_2$ 切开的方式有两种。

因此，我们所考虑的模空间之并不是区间，而是一棵\emph{树}，每当狭缝穿过节点时它就分叉成若干支。当狭缝不继续穿过节点、而是停在节点处、破裂经由该节点处的常值分量发生时，也会给模空间之并的边界带来额外贡献；见图~\ref{fig:ainfty} 中下部的示意图。（为清晰起见，图中把这个构形画在树的一条短边的末端，但实际上并不存在与之相应的某个一维层。）

我们断言：在 $u$ 通过沿 $L_{i-1}$ 或沿 $L_i$ 延伸狭缝而破裂成一对传播圆盘的那些端点之间，仍然存在抵消。面积权重与和乐权重的作用与通常情形完全一样，所以关键要素是引理~\ref{l:tot_mult_invariance}：它断言狭缝穿过节点前后可能出现的各构形的总传播重数相等（模去破裂构形）。这又蕴含：当狭缝朝一个方向延伸时出现的所有破裂构形的传播重数之和（例如图~\ref{fig:ainfty} 左端的三个端点），等于初始圆盘 $u$ 的传播重数；于是，对狭缝朝另一方向延伸作同样的论证，它也等于模空间另一端所有构形的传播重数之和（例如图~\ref{fig:ainfty} 右端的单个端点）。

重复角情形的论证相同，主要差别是：若 $p_i$ 与 $p_{i+1}$ 重合，则狭缝只能沿 $L_i$ 发展；当它一直收缩成一点时，一个对 $\mu^2(p_{i+1},p_i)$ 有贡献的常值圆盘从构形其余部分脱离。引理~\ref{l:tot_mult_invariance} 于是给出：这个无狭缝的破裂构形的传播重数，等于狭缝完全横穿传播圆盘时所有破裂构形的传播重数之和。

这样，关于像有凹角或重复角的传播圆盘的论证，就归结为引理~\ref{l:tot_mult_invariance}，下面证明它。

\subsubsection{引理~\ref{l:tot_mult_invariance} 的证明}\label{sss:proof_lemma_tot}

该引理的证明本质上是组合的，因为我们需要比较狭缝停在节点前与已延伸越过节点两种情形的传播重数。需分别处理三种情形：(1) 狭缝\emph{逆向}穿过节点，即朝离输出标记点 $z_0$ 更远的分量延伸，如图~\ref{fig:ainfty} 左侧；(2) 狭缝\emph{正向}穿过节点，即朝更靠近输出标记点的分量延伸，如图~\ref{fig:ainfty} 右侧；(3) 狭缝穿过一个与输出处的常值圆盘分量相连的节点。每种情形下，所需的传播重数等式都归结为一个纯组合恒等式（引理~\ref{l:ainfty_combi_backwards}–\ref{l:ainfty_combi_outputbounce}）。

\paragraph{情形 1：狭缝逆向穿过节点。}
考虑 $S$ 的某个节点，它映到 $M$ 的节点 $p_v$；$S$ 的一个分量 $D_{in}$ 映到 $M_{e_{in}}$，其边界在 $L_{j_1}$ 与 $L_{j_2}$ 上（$j_1<j_2$），它的输出接在映到 $M_{e_{out}}$ 的分量 $D_{out}$ 的某个输入端上。记 $u$ 在节点附近两个带状端上的次数为 $k_{in}\geq 1$ 与 $k_{out}\geq 0$。设狭缝沿 $L_i$ 从分量 $D_{out}$ 逆向穿过节点伸入 $D_{in}$。由于狭缝来自 $D_{out}$，必有 $i<j_1$ 或 $i>j_2$。当 $i>j_2$（如图~\ref{fig:ainfty} 左侧；$i<j_1$ 时类似），狭缝一旦伸入 $D_{in}$，就把 $p_v$ 附近的局部图像裂成两个传播圆盘：

\begin{itemize}
\item 一个以 $L_{j_2}$ 与 $L_i$ 为边界（$i<j_1$ 时以 $L_i$ 与 $L_{j_1}$ 为边界），从 $M_{e_{out}}$ \emph{逆向}穿过节点传播到 $M_{e_{in}}$，在 $M_{e_{out}}$ 中的输入次数 $1\leq a\leq k_{out}$ 由进入的狭缝在 $D_{out}$ 中的位置决定，在 $M_{e_{in}}$ 中的输出次数 $0\leq b\leq k_{in}-1$ 任意（把狭缝伸入 $D_{in}$ 有 $k_{in}-1$ 种选择）；
\item 另一个以 $L_{j_1}$ 与 $L_i$ 为边界（$i<j_1$ 时以 $L_i$ 与 $L_{j_2}$ 为边界），从 $M_{e_{in}}$ 正向传播到 $M_{e_{out}}$，在 $M_{e_{in}}$ 中的输入次数为 $k_{in}-b$，在 $M_{e_{out}}$ 中的输出次数为 $k_{out}-a$。
\end{itemize}

当 $a=k_{out}$ 时还有另一种可能（对应图~\ref{fig:ainfty} 中下部的图）：狭缝终止于节点 $p_v$，把构形裂成：
\begin{itemize}
\item 一个入射传播圆盘（包含 $u$ 中位于节点 $D_{in}$ 一侧的全部分量，加上 $D_{out}$ 一侧由狭缝界定的两片之一；图~\ref{fig:ainfty} 中下部灰色部分），它从 $D_{in}$ 与 $D_{out}$ 两侧进入节点，并以 $p_v$ 处的常值分量结束；
\item 一个出射传播圆盘（$D_{out}$ 一侧曲线的其余部分），它在 $p_v$ 处有一个输入，边界在 $L_{j_1}$ 与 $L_i$ 上（$i<j_1$ 时为 $L_i$ 与 $L_{j_2}$），局部次数 $k_{out}-a=0$，正如定义~\ref{def:proptree} 对节点处输入所要求的那样。
\end{itemize}

回忆初始构形（局部次数 $k_{in},k_{out}$）在节点 $v$ 处的传播系数 $C^{v;e_{in},e_{out}}_{k_{in},k_{out}}$ 定义为 $t_{e_{in}/v}^{-k_{in}}$ 作为 $t_{e_{out}/v}$ 的幂级数展开式中 $t_{e_{out}/v}^{k_{out}}$ 的系数；而按上述方式插入狭缝后两个节点的传播系数之积为
\[
C^{v;e_{out},e_{in}}_{a,b}\,C^{v;e_{in},e_{out}}_{k_{in}-b,k_{out}-a}.
\]
另一方面，在涉及 $p_v$ 处常值分量的破裂构形（对应 $a=k_{out}$）中，常值分量的各节点对传播重数的贡献为 $K^{v;e_{in},e_{out}}_{k_{in},k_{out}}$，即 $t_{e_{in}/v}^{-k_{in}}t_{e_{out}/v}^{-k_{out}}$ 表成 $t_{e_{in}/v}$ 与 $t_{e_{out}/v}$ 的负幂的线性组合时的常数项。于是，狭缝穿过节点时总传播重数的不变性归结为：

\begin{lemma}\label{l:ainfty_combi_backwards}
给定 $v,e_1,e_2$ 与整数 $k_1\geq 1$、$k_2\geq a\geq 1$，
\begin{equation}\label{eq:ainfty_combi_backwards}
\sum_{b=0}^{k_{1}-1} C^{v;e_{2},e_{1}}_{a,b}\, C^{v;e_{1},e_{2}}_{k_{1}-b,k_{2}-a}
+\delta_{a,k_{2}}\,K^{v;e_{1},e_{2}}_{k_{1},k_{2}} =
C^{v;e_{1},e_{2}}_{k_{1},k_{2}}.
\end{equation}
\end{lemma}

\begin{proof}
记 $t_1=t_{e_1/v}$、$t_2=t_{e_2/v}$。有理函数 $t_1^{-k_1}t_2^{-a}$ 有部分分式分解，可写成单项式 $1,t_1^{-1},\dots,t_1^{-k_1}$ 与 $t_2^{-1},\dots,t_2^{-a}$ 的有限线性组合，即
\[
t_1^{-k_1}t_2^{-a}=K^{v;e_1,e_2}_{k_1,a}+P_1(t_1^{-1})+P_2(t_2^{-1}),
\]
其中 $P_1(t_1^{-1})$ 是 $t_1^{-1}$ 的无常数项多项式（在 $x_1$ 处的极部），$P_2(t_2^{-1})$ 是 $t_2^{-1}$ 的无常数项多项式（在 $x_2$ 处的极部）。另一方面，在 $x_1$ 附近有幂级数展开 $t_2^{-a}=\sum_{b\geq 0} C^{v;e_2,e_1}_{a,b} t_1^b$，从而得到 Laurent 级数
\[
t_1^{-k_1} t_2^{-a}=\sum_{b\geq 0} C^{v;e_2,e_1}_{a,b} t_1^{b-k_1}.
\]
比较 $x_1$ 处的极部（即利用 $P_2$ 在 $x_1$ 附近展成无负幂的 $t_1$ 的幂级数），得
\[
P_1(t_1^{-1})=\sum_{b=0}^{k_1-1} C^{v;e_2,e_1}_{a,b} t_1^{b-k_1}.
\]
再把它代回，利用 $P_1$ 中每个单项式在 $x_2$ 附近的展开 $t_1^{b-k_1}=\sum_{c\geq 0}C^{v;e_1,e_2}_{k_1-b,c}\,t_2^c$，得到 $t_1^{-k_1}t_2^{-a}$ 在 $x_2$ 附近的 Laurent 级数
\[
t_1^{-k_1}t_2^{-a}=P_2(t_2^{-1})+K^{v;e_1,e_2}_{k_1,a}+
\sum_{b=0}^{k_1-1}\sum_{c\geq 0}
C^{v;e_2,e_1}_{a,b} C^{v;e_1,e_2}_{k_1-b,c}\, t_2^c.
\]
另一方面，从 $t_1^{-k_1}=\sum_{d\geq 0} C_{k_1,d}^{v;e_1,e_2}\,t_2^d$ 出发又有
\[
t_1^{-k_1} t_2^{-a}=\sum_{d\geq 0} C^{v;e_1,e_2}_{k_1,d} t_2^{d-a}.
\]
比较两式中 $t_2^{k_2-a}$ 的系数，即得 \eqref{eq:ainfty_combi_backwards}。
\end{proof}

\paragraph{情形 2：狭缝正向穿过节点。}
仍考虑 $S$ 的某个节点，它映到 $M$ 的节点 $p_v$；$S$ 的分量 $D_{in}$（映到 $M_{e_{in}}$，边界在 $L_{j_1},L_{j_2}$ 上，$j_1<j_2$）的输出接在映到 $M_{e_{out}}$ 的分量 $D_{out}$ 的输入端上。再设 $u$ 在 $D_{out}$ 上的限制不是常值映射，并记节点附近两个带状端的次数为 $k_{in}\geq 1$ 与 $k_{out}\geq 0$。设狭缝沿 $L_i$ 从 $D_{in}$ 正向穿过节点伸入 $D_{out}$。由于狭缝来自 $D_{in}$，必有 $j_1<i<j_2$（见图~\ref{fig:ainfty} 右侧）；一旦伸入 $D_{out}$，它就把 $p_v$ 附近的局部图像裂成两个传播圆盘，二者都从 $M_{e_{in}}$ 正向穿过节点到 $M_{e_{out}}$：一个的输入次数为 $1\leq a\leq k_{in}-1$（由狭缝在 $D_{in}$ 中的位置决定），输出次数 $0\leq b\leq k_{out}$（可自由选取，把狭缝伸入 $D_{out}$ 有 $k_{out}+1$ 种方式）；另一个的输入次数为 $k_{in}-a$，输出次数为 $k_{out}-b$。于是总传播重数的不变性归结为：

\begin{lemma}\label{l:ainfty_combi_forward}
给定 $v,e_1,e_2$ 与整数 $1\leq a\leq k_1-1$、$k_2\geq 0$，
\begin{equation}\label{eq:ainfty_combi_forward}
\sum_{b=0}^{k_2} C_{a,b}^{v;e_1,e_2} C_{k_1-a,k_2-b}^{v;e_1,e_2}=C_{k_1,k_2}^{v;e_1,e_2}.
\end{equation}
\end{lemma}

\begin{proof}
把 $t_1^{-k_1}$ 写成 $t_1^{-a}=\sum_{b\geq 0} C_{a,b}^{v;e_1,e_2}\, t_2^b$ 与 $t_1^{a-k_1}=\sum_{d\geq 0} C_{k_1-a,d}^{v;e_1,e_2}\, t_2^d$ 之积，再取所得幂级数中 $t_2^{k_2}$ 的系数即得。
\end{proof}

\paragraph{情形 3：狭缝穿过常值输出分量。}

\begin{figure}[t]
\centering
\begin{tikzpicture}[scale=0.75]
% 左：常值输出分量
\begin{scope}[shift={(-3.6,0)}]
\fill[orange!18] (0,0) circle (1.15);
\draw[thick] (0,0) circle (1.15);
\fill (0,0) circle (0.07) node[above] {\small $p_0=p_v$};
\draw[blue,thick] (0.0,1.15)--(0,0);
\draw[green!50!gray,thick] (-1.15,0)--(0,0);
\draw[red,thick] (0,-1.15)--(0,-0.001);
\node[green!50!gray] at (-1.5,0.3) {\small $L_0$};
\node[red] at (-0.35,-1.5) {\small $L_k$};
\node[blue] at (0.45,1.0) {\small $L_i$};
\node at (0,-1.9) {\small (a) 常值分量有两个输入};
\end{scope}
\draw[->,thick] (-1.4,0)--(-0.5,0);
\begin{scope}[shift={(2.6,0)}]
\fill[orange!18] (0,0) circle (1.15);
\draw[thick] (0,0) circle (1.15);
\fill (0,0) circle (0.07);
\draw[blue,thick] (0.0,1.15)--(0,0.15);
\draw[blue,thick] (0.0,0.15)--(0.55,-0.95);
\draw[green!50!gray,thick] (-1.15,0)--(-0.1,0);
\draw[red,thick] (0,-1.15)--(-0.1,0);
\node[green!50!gray] at (-1.5,0.3) {\small $L_0$};
\node[red] at (-0.35,-1.5) {\small $L_k$};
\node[blue] at (0.75,0.75) {\small $L_i$};
\node at (0,-1.9) {\small (b) 狭缝穿过它};
\end{scope}
\end{tikzpicture}
\caption{狭缝穿过常值输出分量（示意图；原书图 5）。}\label{fig:slitoutput}
\end{figure}

接着考虑狭缝沿 $L_i$ 伸入常值输出分量 $D_{out}$ 的情形（带两个输入、承载输出标记点 $z_0$、映到节点 $p_0=p_v$ 的常值圆盘）。记 $D_1,D_2$ 为 $S$ 中与 $D_{out}$ 相邻的两个分量，$M_{e_1},M_{e_2}$ 为它们所映到的 $M$ 的分量，$k_1,k_2\geq 1$ 为 $u$ 在它们输出带状端上的限制的次数。沿 $L_i$ 在分量 $D_2$ 内延伸并抵达常值输出分量的狭缝，可以进一步伸入 $D_1$，如图~\ref{fig:slitoutput} 所示。这样就把 $p_v$ 附近的局部图像分解为两个传播圆盘：
\begin{itemize}
\item 一个以 $L_i$ 与 $L_j$ 为边界，从 $D_2$ 穿过节点传播到 $D_1$，在 $M_{e_2}$ 中的输入次数 $1\leq a\leq k_2-1$（由狭缝在 $D_2$ 中的位置决定），在 $M_{e_1}$ 中的输出次数 $0\leq b\leq k_1-1$（有 $k_1$ 种可能选择）；
\item 另一个的入射端分别以 $L_0$ 与 $L_i$、以及 $L_i$ 与 $L_k$ 为边界，在 $M_{e_1}$ 与 $M_{e_2}$ 中的次数分别为 $k_1-b$ 与 $k_2-a$，终止于 $p_0=p_v$ 处的常值分量。
\end{itemize}
总传播重数之和不变量于是归结为：

\begin{lemma}\label{l:ainfty_combi_outputbounce}
给定 $v,e_1,e_2$ 与整数 $k_1\geq 1$、$1\leq a\leq k_2-1$，
\begin{equation}\label{eq:ainfty_combi_outputbounce}
\sum_{b=0}^{k_1-1} C_{a,b}^{v;e_2,e_1} K_{k_1-b,k_2-a}^{v;e_1,e_2}=
K_{k_1,k_2}^{v;e_1,e_2}.
\end{equation}
\end{lemma}

\begin{proof}
与引理~\ref{l:ainfty_combi_backwards} 的证明一样，令 $t_1=t_{e_1/v}$、$t_2=t_{e_2/v}$，从部分分式分解 $t_1^{-k_1}t_2^{-a}=K_{k_1,a}^{v;e_1,e_2}+P_1(t_1^{-1})+P_2(t_2^{-1})$ 出发，并回忆
\[
P_1(t_1^{-1})=\sum_{b=0}^{k_1-1} C_{a,b}^{v;e_2,e_1} t_1^{b-k_1}.
\]
乘以 $t_2^{a-k_2}$，得
\begin{equation}\label{eq:ainfty_combi_lemma}
t_1^{-k_1}t_2^{-k_2}=\Bigl(K_{k_1,a}^{v;e_1,e_2}+P_2(t_2^{-1})\Bigr) t_2^{a-k_2}+
\sum_{b=0}^{k_1-1} C_{a,b}^{v;e_2,e_1} t_1^{b-k_1}\,t_2^{a-k_2}.
\end{equation}
这个表达式又可作部分分式分解，写成 $1,t_1^{-1},\dots,t_1^{-k_1},t_2^{-1},\dots,t_2^{-k_2}$ 的线性组合；我们关心其中的常数项。\eqref{eq:ainfty_combi_lemma} 右端第一部分只含 $t_2$ 的负幂，故对常数项无贡献。而 $t_1^{b-k_1}\,t_2^{a-k_2}$ 的部分分式分解中的常数项是 $K_{k_1-b,k_2-a}^{v;e_1,e_2}$；因此 \eqref{eq:ainfty_combi_lemma} 右端部分分式分解中的常数项为
\[
\sum_{b=0}^{k_1-1} C_{a,b}^{v;e_2,e_1}\,K_{b-k_1,a-k_2}^{v;e_1,e_2},
\]
它正是 \eqref{eq:ainfty_combi_outputbounce} 的左端。另一方面，$t_1^{-k_1}t_2^{-k_2}$（\eqref{eq:ainfty_combi_lemma} 左端）的部分分式分解中的常数项，按定义等于 $K_{k_1,k_2}^{v;e_1,e_2}$。比较这两个量即得引理。
\end{proof}

\subsubsection{分叉的传播圆盘}\label{sss:bifurcated}

除带凹角的传播圆盘之外，还有另一类构形会给出传播圆盘的一维模空间：``分叉''圆盘，即在 $M$ 的某个节点 $p_v$ 附近，$S$ 有\emph{两个}入射分量 $D_1,D_2$ 与一个出射分量 $D_{out}$，三者分别映到 $M$ 的不同分量（$M_{e_1}$、$M_{e_2}$、$M_{e_{out}}$，其中 $e_1,e_2,e_{out}$ 是在 $v$ 处相交的三条边）。

若分叉节点附近的出射分量在局部不满射到 $M_{e_{out}}$ 中该节点的邻域（即输出次数 $k_{out}=0$），则这样的分叉圆盘可以直接实现为两个刚性传播圆盘的破裂构形：一个含 $D_1$ 与 $D_2$、终止于 $p_v$ 处的常值分量；另一个从 $p_v$ 处的输入开始、含 $D_{out}$（见图~\ref{fig:ainfty_bifurcated}(c)）。一般地（不论 $k_{out}$ 取值如何），该构形也可以通过把狭缝长入三个分量 $D_1$、$D_2$ 或 $D_{out}$ 中的任一个来变形，其效果是把分叉构形在局部裂成一对真正的传播圆盘。因此，传播圆盘的模空间朝三类方向延伸，对应狭缝在 $S$ 中产生并延伸的三种方式；见图~\ref{fig:ainfty_bifurcated}(a)(b)(d)。（对其中每一种，若 $u$ 在该分量上的次数大于 $1$，还可能有多重可能性。）随着狭缝在 $S$ 的相应分量中展开，它最终要么碰到定义域的边界，要么像带凹角的圆盘那样穿过其他节点、伸入其他分量。这部分与上一小节完全一样，由引理~\ref{l:tot_mult_invariance} 处理；新的要素反倒是：与把分叉节点局部裂成一对刚性传播圆盘的各种方式相关联的组合传播重数之间必须发生的抵消。

\begin{figure}[t]
\centering
\begin{tikzpicture}[scale=0.95]
% (c) 分叉节点：三条边界弧在 p_v 相遇
\begin{scope}[shift={(-4.4,0)}]
  \fill[orange!10,draw=black!55] (0,0) circle (0.95);
  \fill (0,0) circle (1.6pt);
  \draw[green!45!black,thick] (0,0)--(-0.67,0.67);
  \draw[red,thick] (0,0)--(-0.67,-0.67);
  \draw[blue,thick] (0,0)--(0.95,0);
  \node[green!45!black] at (-1.12,1.12) {\scriptsize $L_i$};
  \node[red] at (-1.12,-1.12) {\scriptsize $L_j$};
  \node[blue] at (1.14,0.2) {\scriptsize $L_k$};
  \node at (-0.55,0.62) {\scriptsize $D_1$};
  \node at (-0.55,-0.62) {\scriptsize $D_2$};
  \node at (0.58,0.40) {\scriptsize $D_{out}$};
  \node at (0,-1.45) {\scriptsize (c) 分叉节点};
\end{scope}
\draw[->,thick] (-2.9,0)--(-2.2,0);
% (a) 狭缝长入 D_1
\begin{scope}[shift={(-0.5,0)}]
  \fill[orange!10,draw=black!55] (0,0) circle (0.95);
  \fill (0,0) circle (1.4pt);
  \draw[green!45!black,thick] (0,0)--(-0.67,0.67);
  \draw[red,thick] (0,0)--(-0.67,-0.67);
  \draw[blue,thick] (0,0)--(0.95,0);
  \draw[blue,thick,dashed] (0.05,0.05)--(-0.30,0.62);
  \node at (0,-1.45) {\scriptsize (a) 狭缝长入 $D_1$};
\end{scope}
% (b) 狭缝长入 D_2
\begin{scope}[shift={(3.0,0)}]
  \fill[orange!10,draw=black!55] (0,0) circle (0.95);
  \fill (0,0) circle (1.4pt);
  \draw[green!45!black,thick] (0,0)--(-0.67,0.67);
  \draw[red,thick] (0,0)--(-0.67,-0.67);
  \draw[blue,thick] (0,0)--(0.95,0);
  \draw[green!45!black,thick,dashed] (0.05,-0.05)--(-0.30,-0.62);
  \node at (0,-1.45) {\scriptsize (b) 狭缝长入 $D_2$};
\end{scope}
% (d) 狭缝长入 D_out
\begin{scope}[shift={(6.5,0)}]
  \fill[orange!10,draw=black!55] (0,0) circle (0.95);
  \fill (0,0) circle (1.4pt);
  \draw[green!45!black,thick] (0,0)--(-0.67,0.67);
  \draw[red,thick] (0,0)--(-0.67,-0.67);
  \draw[blue,thick] (0,0)--(0.95,0);
  \draw[red,thick,dashed] (0.08,0.06)--(0.55,0.28);
  \node at (0,-1.45) {\scriptsize (d) 狭缝长入 $D_{out}$};
\end{scope}
\end{tikzpicture}
\caption{把分叉圆盘分解为一对传播圆盘（示意图；原书图 6）。(c) 三条边界弧 $L_i,L_j,L_k$ 在节点 $p_v$ 相遇，形成分叉；(a)(b)(d) 狭缝长入 $D_1$、$D_2$ 或 $D_{out}$ 之一，把分叉构形局部裂成一对刚性传播圆盘。}\label{fig:ainfty_bifurcated}
\end{figure}


记 $D_1,D_2,D_{out}$ 在分叉节点附近的次数为 $k_1\geq 1$、$k_2\geq 1$、$k_{out}\geq 0$。如上所述，若 $k_{out}=0$，则存在一个破裂构形，其中一个刚性传播圆盘含 $D_1$ 与 $D_2$、以 $p_v$ 处的常值分量结束（图~\ref{fig:ainfty_bifurcated}(c)）；与常值分量相邻的节点对该破裂构形的传播重数贡献 $K^{v;e_1,e_2}_{k_1,k_2}$。另一方面，对每个 $0\leq b\leq k_1-1$，都有使狭缝长入 $D_1$ 分量的变形，把局部图像分解为：由 $D_2$（入射到 $p_v$）接上 $D_1$ 的一部分（以次数 $b$ 出射）构成的传播圆盘，以及由 $D_1$ 的其余部分（以次数 $k_1-b$ 入射到 $p_v$）接上 $D_{out}$ 构成的传播圆盘（图~\ref{fig:ainfty_bifurcated}(b)）。类似地，有狭缝在 $D_2$ 分量中的构形：一个传播圆盘由 $D_1$（入射到 $p_v$）接上 $D_2$ 的一部分（以次数 $0\leq a\leq k_2-1$ 出射）构成，另一个由 $D_2$ 的其余部分（以次数 $k_2-a$ 入射到 $p_v$）接上 $D_{out}$ 构成（图~\ref{fig:ainfty_bifurcated}(a)）。最后一种情形是狭缝位于 $D_{out}$ 中：一个传播圆盘由 $D_1$ 接上 $D_{out}$ 的一部分（以次数 $0\leq c\leq k_{out}$ 出射）构成，另一个由 $D_2$ 接上 $D_{out}$ 的其余部分（以次数 $k_{out}-c$ 出射）构成（图~\ref{fig:ainfty_bifurcated}(d)）。于是，比较狭缝位于某个输入圆盘中与位于输出分量 $D_{out}$ 中的构形的传播重数之和，就归结为验证下述恒等式：

\begin{lemma}\label{l:ainfty_tripod}
给定 $G$ 的顶点 $v$ 及其相邻的边 $e_1,e_2,e_3$，以及整数 $k_1,k_2\geq 1$、$k_3\geq 0$，
\begin{equation}\label{eq:ainfty_tripod}
\sum_{a=0}^{k_2-1} C_{k_1,a}^{v;e_1,e_2} C_{k_2-a,k_3}^{v;e_2,e_3}
+\sum_{b=0}^{k_1-1} C_{k_2,b}^{v;e_2,e_1} C_{k_1-b,k_3}^{v;e_1,e_3}
+\delta_{k_3,0}K_{k_1,k_2}^{v;e_1,e_2}=
\sum_{c=0}^{k_3} C_{k_1,c}^{v;e_1,e_3} C_{k_2,k_3-c}^{v;e_2,e_3}.
\end{equation}
\end{lemma}

\begin{proof}
等式来自把 $t_1^{-k_1}t_2^{-k_2}$ 表成 $t_3$ 的幂级数的两种方式。一方面，从 $t_1^{-k_1}=\sum_{c\geq 0} C_{k_1,c}^{v;e_1,e_3} t_3^c$ 与 $t_2^{-k_2}=\sum_{d\geq 0} C_{k_2,d}^{v;e_2,e_3} t_3^d$ 出发，相乘后所得 $t_3$ 的幂级数中 $t_3^{k_3}$ 的系数恰为 \eqref{eq:ainfty_tripod} 的右端。另一方面，可以像引理~\ref{l:ainfty_combi_backwards} 的证明那样，得到部分分式分解
\begin{equation}\label{eq:partialproduct}
t_1^{-k_1}t_2^{-k_2}=K_{k_1,k_2}^{v;e_1,e_2}+P_1(t_1^{-1})+P_2(t_2^{-1})=
K_{k_1,k_2}^{v;e_1,e_2}+\sum_{b=0}^{k_1-1} C_{k_2,b}^{v;e_2,e_1}\,t_1^{b-k_1}
+\sum_{a=0}^{k_2-1} C_{k_1,a}^{v;e_1,e_2}\,t_2^{a-k_2}.
\end{equation}
代入 $t_1^{b-k_1}=\sum_{d\geq 0} C_{k_1-b,d}^{v;e_1,e_3}\, t_3^d$ 与 $t_2^{a-k_2}=\sum_{d\geq 0} C_{k_2-a,d}^{v;e_2,e_3}\, t_3^d$，所得 $t_3$ 的幂级数中 $t_3^{k_3}$ 的系数正是 \eqref{eq:ainfty_tripod} 的左端。
\end{proof}

至此逐情形分析完成，定理~\ref{thm:ainfty} 证毕。
% [Tutor Note: 点拨] 三个组合引理的共同点是：$t_{e_1/v}^{-k_1}t_{e_2/v}^{-k_2}$ 这一有理函数在 $\PP^1$ 上有两种等价写法——在 $x_2$ 处的幂级数展开与在 $x_1$ 处的幂级数展开，二者由部分分式分解（即 $H^1(\PP^1,\cO)=0$）联系。传播系数正是这两组展开的系数，所以传播重数在狭缝穿过节点时守恒，等价于同一个有理函数换一种展开方式。依据：原文 Lemmas 3.12–3.15 及本工程 03_术语表。
